\chapter{Technical Risk Assessment}

EASA\footnote{https://www.easa.europa.eu/en/domains/drones-air-mobility/operating-drone/specific-category-civil-drones, accessed 23.04.2026, accessed 23.04.2026} category: Specific

EU\footnote{https://eur-lex.europa.eu/legal-content/EN/TXT/?uri=CELEX\%3A32020R1058, accessed 23.04.2026}

\section{Specific Operation Risk Assessment (SORA)}
% SORA Classification guide
\footnote{http://jarus-rpas.org/wp-content/uploads/2024/06/SORA-v2.5-Main-Body-Release-JAR\_doc\_25.pdf, to be added in the ref list!}

\subsection{Documentation of Proposed Operations}
Compile operational, technical, and organisational information. This may include:
\begin{itemize}
 \item Various maps, figures, diagrams and other information detailing the operational volume, ground risk buffers, adjacent area, and adjacent airspace.
 \item Information on the operational, technical, and organisational elements of: A. The operation and functions during flight, including intended flight profiles, states, and modes that provide safety throughout the nominal, contingency, and emergency phases of flight, B. Any ground and air risk mitigations (strategic and tactical) used to reduce the intrinsic ground risk or initial air risk.
 \item A description of the contingency volume and ground risk buffers, and how they were determined. 
 \item The applicant may use Annex A, Section A.3 to assist in understanding the type of data that needs to be presented and any other information that supports the risk assessment to the authority.
\end{itemize}


\subsection{Intrinsic Ground Risk Class (iGRC)}
JARUS Annex F \footnote{http://jarus-rpas.org/wp-content/uploads/2024/06/SORA-v2.5-Annex-F-Release.JAR\_doc\_29pdf.pdf}

JARUS other source \footnote{http://jarus-rpas.org/wp-content/uploads/2023/07/jar\_doc\_6\_sora\_sts\_01\_edition1.1.pdf}

Maximum characteristic dimension:

Maximum velocity: 200+ m/s

iGRC footprint: Identify the flight geography, Calculate the contingency volume, Calculate the initial ground risk buffer (1 to 1 principle from other source). \textcolor{red}{Need an FL estimation!}


Maximum population density: Controlled Ground Area

Resultant iGRC: between 3 and 4 (estimate - 3 for drones flying up to 200m/s).


\subsection{Final Ground Risk Class (GRC) (optional)}
Only if our iGRC is too high.

\subsection{Determination of the initial Air Risk Class (ARC)}
reserved or restricted airspaces - atypical airspaces
Atypical airspace: ARC-a

\subsection{Application of strategic mitigations to determine residual ARC (optional)}

\subsection{Tactical Mitigation Performance Requirement (TMPR) and Robustness Levels}
No TMPR requirement for ARC-a

From Annex D \footnote{http://jarus-rpas.org/wp-content/uploads/2024/06/SORA-Annex-D-v1.0.pdf}, section 5: "No System Risk Ratio guidance; although operator/applicant may still need to show some form of mitigation as deemed necessary by the CAA"

Utilise Annex A \footnote{http://jarus-rpas.org/wp-content/uploads/2024/06/SORA-v2.5-Annex-A-Release.JAR\_doc\_26-pdf.pdf}, Section A.3 for further guidance on presenting the data supplementing the risk assessment to the authority.

\subsection{Specific Assurance and Integrity Levels (SAIL) determination}
Flying over military area only gives us ARC-a and GRC of 3 or 4, which gives us SAIL II or SAIL III - most likely SAIL III.


\subsection{Determination of Containment Requirements}
Size of adjacent area: 35km
Containment Requirements: Medium/High

Criterion 1 (Operational Volume Containment): (Qualitative) No probable1 single failure of the UAS or any external system supporting the operation shall lead to operation outside of the operation volume. OR (Quantitative) The probability of the failure condition “UA leaving the operational volume” shall be less than 10-3/Flight Hour (FH).

Criterion 2 (End of Flight upon exit of the operational volume): When the UA leaves the operational volume, an immediate end of the flight must be initiated through a combination of procedures/processes and/or available technical means.

Criterion 3 (Definition of the final ground risk buffer): Ground risk buffer must consider the following points below:
\begin{itemize}
    \item Probable single failures (including the projection of high energy parts such as rotors and propellers) which would lead to an operation outside of the operational volume,
    \item Meteorological conditions (e.g., maximum sustained wind),
    \item UAS latencies (e.g., latencies that affect the timely manoeuvrability of the UA),
    \item UA behaviour when activating a technical containment measure, UA performance.
\end{itemize}

Criterion 4 (Ground risk buffer containment): No single failure of the UAS or any external system supporting the operation shall lead to operation outside of he ground risk buffer. Software (SW) and Airborne Electronic Hardware (AEH) whose development error(s) could directly lead to operations outside of the ground risk buffer shall be developed to an industry standard or methodology recognized as adequate by the competent authority.

Criterion 5: 




Criterion 19: The applicant declares1
that the
required level of integrity has been
achieved.

Criterion 20: The applicant conducts a human factors
evaluation of the UAS to determine if
the HMI is appropriate for the mission.
The HMI evaluation is based on
inspection or analyses. The adequacy of
the result of the HMI evaluation is
declared.

Criterion 23: Environmental conditions for safe operations are defined and reflected in the flight manual or equivalent document.1



criterion 24: The UAS is designed to perform as
intended in the environmental
conditions defined and reflected in the
flight manual or equivalent document.


\subsection{Identification of Operational Safety Objectives (OSO)}

\subsection{Comprehensive Safety Portfolio (CSP)}